\documentclass[11pt,a4paper]{article}
\usepackage[margin=1in]{geometry}
\usepackage[T1]{fontenc}
\usepackage[utf8]{inputenc}
\usepackage{amsmath,amssymb,booktabs,graphicx,array,hyperref}
\hypersetup{colorlinks=true,linkcolor=blue,citecolor=blue,urlcolor=blue}
\title{\textbf{SACG: Spatially Adaptive Consistency Gating for Refinement of Frozen Prompt-Based Metric Depth Estimates}}
\author{Author list to be finalized}
\date{Working draft --- October 2026}
\begin{document}
\maketitle
\begin{abstract}
Prompt-conditioned metric depth estimators combine image priors with low-resolution range measurements but can retain local depth errors. We study SACG, a lightweight refinement module trained while Prompt Depth Anything (PromptDA) remains entirely frozen. SACG uses RGB--depth gradient agreement and a learned spatial field derived from the validity of LiDAR samples to predict a gate, which modulates a depth residual estimated from RGB, coarse depth and frozen DPT features. Our selected variant is trained with valid-pixel L1 and a delayed gate penalty, without the boundary-weighted term of an earlier version. On 5,599 ARKitScenes validation frames (one training seed), it reduces AbsRel from 0.011505 to 0.011298 and L1 from 0.013479 to 0.013227\,m. It narrowly outperforms a separately trained residual CNN on these two metrics, but not on RMSE, RGB-edge-based boundary error or SSIM. We report these observations with their limitations and identify the controlled additional experiments necessary to evaluate architectural advantage.
\end{abstract}
\section{Introduction}
Dense metric depth is central to robotics and scene understanding. RGB alone cannot resolve scale, while mobile range sensors supply low-resolution or sparse depth. Image-guided depth completion combines these signals \cite{ma2018sparse,park2020nlspn}.

Pretrained visual representations, including DINOv2 \cite{oquab2024dinov2} and Depth Anything \cite{yang2024depthanything}, enable capable dense depth predictors. PromptDA \cite{lin2025promptda} accepts metric depth prompts and provides accurate dense estimates. Yet downstream images retain localized residual errors. We investigate a small learnable postprocessing module that improves these predictions while keeping PromptDA completely frozen.

Our spatially adaptive consistency gating (SACG) module predicts a residual from RGB, coarse depth and frozen DPT decoder features, then modulates it using a learned gate. The gate uses RGB--depth gradient agreement and a spatial field derived from the validity pattern of the low-resolution LiDAR measurements. The refinement CNN itself does not ingest raw LiDAR depth values.

In the available ARKitScenes seed-10 ablation, the model trained without the optional boundary-weighted loss slightly outperforms both frozen PromptDA and a separately trained simple residual CNN on L1 and AbsRel. We therefore select it as the working method, but do not infer a general architectural advantage from this single-seed comparison.

Contributions of this working manuscript are (1) a code-grounded specification of frozen-PromptDA residual refinement with a consistency-conditioned spatial gate, (2) a transparent comparison to an independently trained plain CNN and component ablations, and (3) a measured report of metric-specific improvements and limitations. Claims about zero-shot generalization, calibrated confidence and state-of-the-art accuracy are out of scope.

\section{Related Work}
\label{sec:related}
The proposed method connects three areas: learning from incomplete depth measurements, metric depth estimation with pretrained visual representations, and spatially selective refinement of an existing prediction. We distinguish these areas because SACG does not train a depth estimator from scratch. It refines an already strong, prompt-conditioned foundation model while keeping its backbone and decoder frozen.

\subsection{Depth Completion From Sparse Measurements}
The depth-completion task aims to reconstruct a dense metric field from sparse or low-resolution range measurements, often with guidance from a registered RGB image. Sparse-to-Dense~\cite{ma2018sparse} established a trainable encoder--decoder formulation for the joint processing of RGB and depth observations. The approach demonstrated that even relatively few metric samples can constrain dense inference when interpreted together with image structure. A key distinction from our setting is that a conventional completion network learns the entire depth mapping, whereas our adapter modifies only the residual of a pretrained metric estimator.

Convolutional spatial propagation networks (CSPN)~\cite{cheng2018cspn} treat dense depth refinement as repeated diffusion according to learned local affinities. Importantly, their propagation formulation can incorporate known sparse samples as constraints. NLSPN~\cite{park2020nlspn} extends propagation to non-local neighborhoods using learned confidence and affinities, showing the importance of spatially adaptive behavior when observed depth is unevenly distributed. These methods establish that learned spatial weighting and confidence-informed correction are well known in depth completion. Therefore, the novelty of SACG cannot be the mere use of a gate or the notion of confidence-guided refinement.

More recently, CompletionFormer~\cite{zhang2023completionformer} combines convolutional operators and transformers to leverage local details and non-local interactions. Such architectures are designed to propagate sparse geometry over large regions. Direct performance comparisons between these systems and a PromptDA adapter require consistent training data, sensor sparsity, resolution, and evaluation protocol. We do not infer superiority over published completion benchmarks from the paired validation results reported in this manuscript.

\subsection{Depth Foundation Models and Metric Depth Prompting}
Strong visual representation learning has changed how depth prediction systems are built. DINOv2~\cite{oquab2024dinov2} provides general visual features, whereas Depth Anything~\cite{yang2024depthanything} and Depth Anything V2~\cite{yang2024depthanythingv2} demonstrate the benefits of large-scale depth pretraining. These methods improve single-image geometry estimation but may still require metric anchoring when absolute depth is needed.

Prompt Depth Anything (PromptDA)~\cite{lin2025promptda} integrates low-cost LiDAR measurements into the depth foundation model's decoder and produces metric depth at substantially higher resolution than the input prompt. Its pretrained representation already captures rich visual and metric information. This creates a distinct adaptation setting: instead of discarding the pretrained prediction, we can inspect whether a small learned correction improves local residual errors. Our work uses the pretrained PromptDA ViT-L inference pathway without modifying its prompt-fusion layers, visual backbone, or decoder.

\subsection{Residual Prediction, Gating, and Confidence}
Given a coarse prediction $D_0$, residual refinement constructs an estimated correction $R$ and forms $D_0+R$. A spatially gated variant uses $D_0+G\odot R$, where $G$ is a learned coefficient at each location. The second formulation offers an explicit spatial modulation mechanism, but a sufficiently expressive residual CNN can potentially learn similar behavior without an independently parameterized gate.

The proposed SACG architecture separates the two roles. A residual CNN uses RGB, coarse depth, and intermediate frozen DPT features, while the gate network receives a normalized RGB--depth gradient-agreement map, coarse depth, and a learned field constructed from the locations of valid LiDAR measurements. This differs from direct propagation: SACG neither iteratively propagates measured depths nor explicitly enforces equality at sensor-observed pixels. The gate is also not calibrated predictive uncertainty, since no uncertainty labels, distributional likelihood, or calibration objective are used.

The distinction between LiDAR sample values and their spatial support is important. In the implemented simple CNN baseline, resized numerical LiDAR observations and the validity mask are direct inputs. SACG's residual CNN does not receive the numerical depth prompt directly; the frozen PromptDA receives those values upstream, and the SACG gate sees only a validity-derived field. The two trained heads consequently differ in feature availability in addition to gating. Their performance comparison is useful as a system-level control, but alone cannot isolate the gate's causal benefit.

\subsection{Datasets and Evaluation Perspective}
ARKitScenes~\cite{baruch2021arkitscenes} provides mobile RGB-D sequences and high-quality depth supervision relevant to the proposed use of low-cost metric prompts. ScanNet++~\cite{yeshwanth2023scannetpp} combines high-fidelity indoor scene geometry with commodity-sensor image and depth captures. Results depend on the particular scene splits, sensor alignment, valid-depth conventions, and reference-depth construction. Accordingly, we retain common frame manifests for paired comparisons and distinguish local dataset exports from official published benchmark settings.

The precise research question studied here is whether consistency-conditioned spatial modulation offers a repeatable benefit for supervised downstream refinement of an otherwise frozen prompt-conditioned depth foundation model. The existing experiments support only a preliminary, metric-specific answer; a more controlled study is needed to separate gate design from feature inputs and training-objective choices.

\section{Proposed Method}
\label{sec:method}
\subsection{Problem Setting and Design Rationale}
Let $I\in[0,1]^{3\times H\times W}$ be a registered RGB observation and $S\in\mathbb R_{\ge0}^{1\times h\times w}$ a lower-resolution metric depth prompt. We denote the available dense ground truth by $D^\star$, with an accompanying valid set $\Omega$. The first stage is an off-the-shelf pretrained PromptDA model $P_\theta$:
\begin{equation}
D_0=P_\theta(I,S), \qquad \theta\ \text{fixed}.
\label{eq:base}
\end{equation}
In the implementation, the backbone is explicitly placed in evaluation mode, all its parameters have \texttt{requires\_grad=False}, and its forward pass runs under \texttt{torch.no\_grad()}. In addition to $D_0$, a forward hook at \texttt{depth\_head.scratch.output\_conv1} exposes an intermediate depth-decoder feature $F_{\rm DPT}$. We use a separate lightweight trainable module to produce a revised prediction, rather than changing PromptDA's visual encoder or metric-prompt pathway.

Refining a strong frozen model motivates two complementary functions. A \emph{correction generator} should estimate the direction and magnitude of a potentially useful metric-depth adjustment. A \emph{correction gate} should determine how much of that adjustment is applied at each pixel. Their implemented combination is
\begin{equation}
\widehat D=D_0+G\odot R,
\label{eq:refine}
\end{equation}
where $R\in\mathbb R^{1\times H\times W}$ is the proposed residual and $G\in(0,1)^{1\times H\times W}$ is the spatial gate. The formulation does not force every residual to have the same magnitude. Nor does it, by itself, guarantee that a gated model is better than an ungated residual CNN; we test that empirical question separately.

\subsection{Gradient Agreement as a Contextual Cue}
Depth discontinuities often correspond to image edges, but visual texture, lighting changes, and occlusion can break this correspondence. We compute a deliberately simple image--depth gradient-agreement cue rather than inferring geometric confidence from RGB appearance alone. For a single-channel field $X$, the Sobel magnitude and normalization are
\begin{align}
e(X)&=\sqrt{(K_x*X)^2+(K_y*X)^2+\epsilon},\\
E(X)&=\frac{e(X)}{\max_{p}e(X)_p},
\label{eq:edge}
\end{align}
where $K_x$ and $K_y$ are standard $3\times3$ Sobel filters and $\epsilon$ is a small stabilizer. The maximum is taken separately over the spatial dimensions for each image. The RGB tensor is converted into a grayscale image by averaging its three channels; coarse depth is repeated along the channel dimension before the shared routine, which yields the same mean-channel depth values.

We define the gradient-consistency cue as
\begin{equation}
C=\operatorname{clip}\bigl(1-|E(I)-E(D_0)|,\ 0,\ 1\bigr).
\label{eq:cgrad}
\end{equation}
Large $C$ denotes similar \emph{normalized edge magnitudes}, not necessarily matching edge orientations or physically correct depth. A large score may arise in textureless regions where both signals have weak gradients. Conversely, a shadow boundary may have a strong appearance response and no corresponding depth discontinuity. Consequently, the cue enters a learnable gating network together with other context; it should not be interpreted as a calibrated confidence probability.

\subsection{Spatial Support of LiDAR Measurements}
The sparse metric measurements play two roles in the full system. Their \emph{numerical values} enter the pretrained PromptDA forward pass in Eq.~\eqref{eq:base}. Separately, their \emph{availability pattern} contributes spatial context to the adapter. Specifically, we construct a low-resolution binary validity mask
\begin{equation}M_{ij}=\mathbf{1}\{S_{ij}>0\}.\end{equation}
The selected configuration computes an unnormalized field $U=K_\psi*M$, where $K_\psi$ is a single-channel $31\times31$ learnable convolution initialized with a two-dimensional Gaussian-like kernel. Within each image, the map is min--max normalized and resized:
\begin{align}
\overline U &= \frac{U-\min U}{\max(\max U-\min U,10^{-6})},\\
L&=\operatorname{BilinearResize}\bigl(\operatorname{clip}(\overline U,0,1)\bigr).
\label{eq:lidar}
\end{align}
The convolution operates on the native prompt grid before the field is resized to $H\times W$. It is not equivalent to interpolating the metric depth values or explicitly finding nearest three-dimensional neighbors. The resulting field can respond to local sample density or arrangement; because its filter is learned, it is not necessarily monotonic with Euclidean distance to the nearest measurement.

For comparison, the fixed-field implementation computes a distance transform $d(p,M)$ and $L_{\rm fixed}(p)=\exp[-d(p,M)/\sigma]$ before resizing, with nominal $\sigma=50$ in prompt-grid pixels. Another ablation replaces $L$ by zeros. Importantly, these interventions affect only the adapter gate: they do not remove $S$ from the pretrained PromptDA input. Thus, an experiment named \emph{no LiDAR field} does not represent an image-only model.

\subsection{Separate Gate and Residual Networks}
The gate network receives three spatially aligned single-channel inputs:
\begin{equation}
G=\sigma\left(g_\omega([C,L,D_0])\right).
\label{eq:gate}
\end{equation}
The implementation of $g_\omega$ uses a $3\times3$ convolution with 16 channels, batch normalization, ReLU, and a second $3\times3$ convolution to one output channel. The sigmoid $\sigma$ bounds gate values between zero and one. Both the coarse depth itself and the two guidance fields can influence gate placement and intensity.

The correction branch receives an alternative set of information:
\begin{equation}
R=f_\phi\left([I,D_0,\widetilde F_{\rm DPT}]\right),
\label{eq:residual}
\end{equation}
where $\widetilde F_{\rm DPT}$ is a bilinearly resized frozen decoder feature map. The residual CNN comprises three convolution--batch-normalization--ReLU blocks with 32 hidden channels and a final $3\times3$ convolution producing one metric-depth residual channel. Its output is \emph{not} restricted to be positive, allowing either positive or negative corrections through Eq.~\eqref{eq:refine}. In the current code, the output itself is not clamped to positive depth and is not constrained to equal LiDAR measurements at observed locations.

\begin{figure*}[t]
\centering
\begin{tikzpicture}[>=Latex, font=\small, node distance=8mm,
box/.style={draw,rounded corners=2pt,align=center,minimum height=8mm,inner sep=5pt},
arrow/.style={-Latex,thick}]
\node[box,fill=blue!6] (input) {RGB image $I$\\LiDAR prompt $S$};
\node[box,fill=blue!6,right=15mm of input] (pd) {Pretrained PromptDA\\\textbf{frozen}};
\node[box,fill=blue!6,right=16mm of pd] (coarse) {Coarse depth $D_0$\\DPT feature $F_{\rm DPT}$};
\node[box,fill=orange!9,below=16mm of input] (signals) {RGB--depth gradients $C$\\LiDAR validity field $L$};
\node[box,fill=orange!9,right=24mm of signals] (gate) {Gate CNN $g_\omega$\\$G=\sigma(g_\omega([C,L,D_0]))$};
\node[box,fill=green!8,right=23mm of gate] (residual) {Residual CNN $f_\phi$\\$R=f_\phi([I,D_0,\widetilde F_{\rm DPT}])$};
\node[box,fill=green!8,above=11mm of residual] (output) {Refined depth\\$\widehat D=D_0+G\odot R$};
\draw[arrow] (input)--(pd);
\draw[arrow] (pd)--(coarse);
\draw[arrow] (input)--(signals);
\draw[arrow] (signals)--(gate);
\draw[arrow] (coarse.south) |- (gate.north);
\draw[arrow] (coarse.south east) |- (residual.north);
\draw[arrow] (gate)--(residual);
\draw[arrow] (residual)--(output);
\end{tikzpicture}
\caption{System-level depiction of SACG. The frozen pretrained estimator receives numerical metric LiDAR values, while a learned convolution of the validity mask informs the gating path. The residual branch uses frozen DPT features. The two outputs are combined as an additive gated correction; the arrow from gate to residual represents the final multiplication, not conditioning of the residual CNN itself.}
\label{fig:architecture}
\end{figure*}

\subsection{Supervised Objective and Model Selection}
At valid target pixels $\Omega=\{p:D^\star_p>0.01\,{\rm m},\ D^\star_p\ \text{finite}\}$, our basic supervised term is
\begin{equation}
\loss_{\rm L1}=\frac{1}{|\Omega|}\sum_{p\in\Omega}|\widehat D_p-D^\star_p|.
\label{eq:l1}
\end{equation}
A gate penalty discourages large gate activations where the heuristic gradient agreement is high:
\begin{equation}
\loss_{\rm gate}=\frac{1}{HW}\sum_p G_p\,\operatorname{stopgrad}(C_p).
\label{eq:lgate}
\end{equation}
The cue is detached in the regularizer. The gate loss is \emph{not} a likelihood term and does not train a calibrated confidence estimate. The implementation activates this penalty after a warmup of 5,000 optimizer steps:
\begin{equation}
\loss_{\rm selected}(t)=\loss_{\rm L1}
+\mathbf{1}[t\ge5000]\lambda_g\loss_{\rm gate},
\qquad\lambda_g=0.1.
\label{eq:objective}
\end{equation}

An earlier SACG configuration additionally included an RGB-edge-weighted absolute error:
\begin{equation}
\loss_{\rm boundary}
=\frac{1}{HW}\sum_p
e^{10E(I)_p}|\widehat D_p-D^\star_p|\mathbf{1}[p\in\Omega],
\label{eq:lboundary}
\end{equation}
with $\lambda_b=0.5$. Its image-wide denominator differs from the valid-pixel normalization of Eq.~\eqref{eq:l1}. The exponential term can produce disproportionate weights at strong RGB edges, which may be geometric or textural. We found that omitting this term improved several validation metrics in one ARKitScenes experiment. We therefore use Eq.~\eqref{eq:objective}, with \texttt{use\_boundary\_loss=False}, as the selected model. This is a data-driven objective choice and should not be conflated with removing gradient agreement from the network.

\subsection{Plain CNN Control and Attribution}
To distinguish gains from simply adding trainable parameters, the comparison includes a separate \emph{plain residual CNN}. It receives six channels: RGB ($3$), frozen PromptDA depth ($1$), nearest-resized numerical LiDAR depth ($1$), and nearest-resized validity ($1$). It consists of three convolution blocks of 56 channels and a final residual prediction layer, with zero-initialized final convolution weights so that its initial output equals the pretrained estimate. All PromptDA parameters remain frozen.

This control differs from SACG in more than the presence of a gate: it accesses numerical LiDAR observations directly, and it does not use frozen intermediate decoder features. Accordingly, a difference between the two does not by itself measure the isolated effect of gating. The decisive follow-up is to fix the input feature set and objective, then compare a learned $G$ with a constant gate and with a parameter-matched ungated residual. This limitation shapes our interpretation of the results.

\section{Experimental Protocol}
\label{sec:experiments}
\subsection{Datasets and Evaluation Splits}
Our primary study uses ARKitScenes~\cite{baruch2021arkitscenes}, which pairs mobile RGB imagery and low-resolution active-depth observations with dense reference geometry. We evaluate the stored ARKitScenes runs using 5,599 validation frames. RGB and target depth are processed at $756\times1008$ resolution, whereas the low-resolution prompt remains in its original grid so that PromptDA handles the scale discrepancy. The custom loader applies orientation adjustments based on the recorded camera direction, converts depth from millimeters to meters, and uses nearest-neighbor interpolation for dense targets to avoid synthesizing intermediate depth values through bilinear mixing. Validity is defined by finite target depth above $0.01$\,m.

We also inspect an existing locally prepared ScanNet++~\cite{yeshwanth2023scannetpp} dataset export with RGB images, commodity-sensor depth prompts, and rendered reference depth. Its training schedule uses a scene-balanced subset capped at 20,000 frames, and comparisons sample 2,000 validation frames. The export's scenes are separated using local \texttt{train.txt} and \texttt{val.txt} split files; the loader checks that scene identifiers do not overlap. RGB masks can remove invalid rendered target regions and corresponding low-resolution prompt locations. This protocol is tailored to available hardware and is \emph{not claimed to reproduce} the original PromptDA paper's benchmark setup.

\subsection{Frozen Backbone and Optimization}
We use pretrained PromptDA with the ViT-L encoder and \texttt{legacy} DPT decoder configuration. The official pretrained model is loaded for both the baseline and adapter systems. The trainer asserts that all PromptDA parameters are frozen and constructs an AdamW optimizer only over the added refinement module. The implemented learning rate is $10^{-4}$ with weight decay $10^{-4}$; a cosine annealing scheduler spans 10,000 optimizer steps. We use batch size two and gradient-norm clipping at 1.0. The random seed is applied to Python, NumPy, and PyTorch; the code also enables deterministic cuDNN behavior. Train-time augmentation contains horizontal flips, while validation uses a consistent valid-depth mask without random augmentation.

The full SACG objective combines valid-pixel L1, gate regularization beginning at step 5,000, and an edge-weighted boundary term. The selected configuration retains the same architecture, learnable LiDAR field, optimization parameters, and delayed gate regularizer, but disables the boundary term. We use the checkpoint with the lowest validation L1. Since validation influenced both checkpoint selection and the choice to drop the boundary term, this validation split cannot be described as an independent final test set.

\subsection{Compared Systems}
The \emph{frozen baseline} is official PromptDA inference on RGB and the low-resolution depth prompt, without any downstream trainable module. The \emph{plain residual CNN} is a separate, supervised head over RGB, coarse depth, resized numerical LiDAR samples, and their validity mask; it is trained with valid-pixel L1 while PromptDA stays frozen. The \emph{full SACG} ablation uses the originally proposed boundary-weighted objective. Additional variants replace the learnable validity field with a fixed distance-transform field or remove the gate field entirely. These latter ablations still use full-SACG boundary supervision in the available committed results, so they do not isolate the field within the selected no-boundary objective.

\begin{table}[t]
\caption{Main implementation and evaluation settings. Different training seeds are not automatically equivalent to independent data splits.}
\label{tab:protocol}
\centering\small
\begin{tabular}{@{}p{0.41\columnwidth}p{0.49\columnwidth}@{}}\toprule
Setting & Value\\ \midrule
Frozen backbone & PromptDA ViT-L, legacy DPT\\
Dense image/target size & $756\times1008$\\
Optimizer & AdamW\\
Learning rate/decay & $10^{-4}$ / $10^{-4}$\\
Train batch size & 2\\
Optimization budget & 10,000 steps\\
Schedule & Cosine annealing\\
Gradient clipping & Norm 1.0\\
Gate coefficient & 0.1 from step 5,000\\
Selected boundary weight & 0 (disabled)\\
Selection metric & Validation L1\\
ARKitScenes validation & 5,599 frames\\
ScanNet++ local evaluation & 2,000 frames/run\\
Aggregation & Mean of per-image metrics\\ \bottomrule
\end{tabular}
\end{table}

\subsection{Evaluation Measures}
We compute each metric on valid reference pixels and then report the arithmetic mean of the per-image values. For prediction $\widehat D$ and reference $D^\star$, the per-image quantities are
\begin{align}
\mathrm{L1}&=\frac1{|\Omega|}\sum_{p\in\Omega}|\widehat D_p-D^\star_p|,\\
\mathrm{AbsRel}&=\frac1{|\Omega|}\sum_{p\in\Omega}
\frac{|\widehat D_p-D^\star_p|}{D^\star_p},\\
\mathrm{RMSE}&=\sqrt{\frac1{|\Omega|}\sum_{p\in\Omega}
(\widehat D_p-D^\star_p)^2}.
\end{align}
We additionally use standard depth-threshold accuracies $\delta_1,\delta_2,\delta_3$, and a stricter $\delta_{0.5}$ implemented by the project. The RGB-edge-conditioned BoundaryAbsRel restricts the relative error to pixels inside a dilated Sobel-RGB edge mask (threshold 0.1, dilation kernel 5). It measures errors near \emph{image} discontinuities, not necessarily annotated depth discontinuities. Auxiliary SSIM is calculated after setting invalid reference pixels to zero; it is not an official masked-depth structural-similarity benchmark.

\subsection{Comparability and Stored Artifacts}
For each compared pair, the inference script loads an independent frozen PromptDA baseline and a wrapped PromptDA-plus-refiner model, then asserts numerical agreement between their coarse outputs. It saves per-frame metrics, aggregate metrics, input configuration, identifiers, and a validation manifest hash. In the ARKitScenes seed-10 results, the compared variants share an identical 5,599-frame evaluation manifest, enabling paired descriptive comparison. Their shared identifier manifest has SHA-256 prefix \texttt{5b8a182fd4fba0d}. This does not guarantee statistical significance, nor does it eliminate differences in the data made available to the refinement modules.

The currently committed three-seed full-SACG experiments are not trials of the selected no-boundary objective. On ScanNet++, the \texttt{data\_seed} was changed with the training seed, so differences across runs also reflect different evaluated frame subsets. The resulting variation should not be interpreted as pure training-seed variance. We explicitly avoid combining these runs into an uncertainty estimate for the selected model.

\section{Results}
\subsection{ARKitScenes: selection of the main variant}
\begin{table}[t]\centering
\caption{Committed ARKitScenes validation, seed 10, identical 5,599 frames. Lower is better for all four metrics.}
\label{tab:arkit}
\resizebox{\linewidth}{!}{\begin{tabular}{lcccc}\toprule
Method & AbsRel & L1 (m) & RMSE (m) & BoundaryAbsRel \\\midrule
Frozen PromptDA & 0.011505 & 0.013479 & 0.031950 & 0.020598 \\
Simple residual CNN & 0.011321 & 0.013257 & \textbf{0.031793} & \textbf{0.020335} \\
SACG with boundary loss & 0.011358 & 0.013305 & 0.031846 & 0.020446 \\
\textbf{SACG without boundary loss} & \textbf{0.011298} & \textbf{0.013227} & 0.031803 & 0.020382 \\\bottomrule
\end{tabular}}\end{table}
The selected no-boundary variant improves AbsRel by 1.79\% and L1 by 1.87\% over frozen PromptDA, with small additional improvements over the simple CNN. However, the CNN has lower RMSE and RGB-edge-based boundary error. Thus these are single-seed, metric-specific gains, not evidence of overall superiority.
\subsection{Component ablations}
\begin{table}[t]\centering\caption{ARKitScenes seed-10 SACG configurations; field ablations retain boundary loss.}\begin{tabular}{lccc}\toprule
Variant & AbsRel & L1 (m) & BoundaryAbsRel \\\midrule
Full SACG with boundary & 0.011358 & 0.013305 & 0.020446 \\
No LiDAR gate field & 0.011396 & 0.013351 & 0.020497 \\
Fixed LiDAR gate field & 0.011420 & 0.013385 & 0.020480 \\
No boundary loss & \textbf{0.011298} & \textbf{0.013227} & \textbf{0.020382} \\\bottomrule
\end{tabular}\end{table}
The selected loss changes only the boundary term. The field ablations are instead comparisons under the earlier full loss and therefore do not isolate the field's necessity under the selected objective.
\subsection{Existing multi-seed results belong to the earlier full objective}
\begin{table}[t]\centering\caption{ARKitScenes results for SACG \emph{with boundary loss}, not the selected model.}\begin{tabular}{rccc}\toprule
Seed & AbsRel & L1 (m) & RMSE (m) \\\midrule
10 & 0.011358 & 0.013305 & 0.031846 \\
26 & 0.011399 & 0.013365 & 0.031891 \\
42 & 0.011376 & 0.013315 & 0.031876 \\\bottomrule
\end{tabular}\end{table}
\begin{table}[t]\centering\caption{ScanNet++ local-export results for full SACG \emph{with boundary loss}, 2,000 frames per seed; not the selected model.}\begin{tabular}{rcccc}\toprule
Seed & Baseline AbsRel & Full SACG AbsRel & Baseline L1 & Full SACG L1 \\\midrule
10 & 0.029730 & 0.029613 & 0.048934 & 0.048479 \\
26 & 0.028579 & 0.028510 & 0.048643 & 0.048400 \\
42 & 0.029356 & 0.029263 & 0.048766 & 0.048388 \\\bottomrule
\end{tabular}\end{table}
Stored full-objective models improve AbsRel and L1 against frozen PromptDA on each evaluated subset. The ScanNet++ subset varies by seed, and no committed cross-dataset result was found for selected no-boundary SACG. No new checkpoint validation, paired confidence intervals or held-out tests were executed for this draft.
\section{Discussion and Conclusion}

Removing the exponential RGB-edge boundary loss coincides with lower AbsRel and L1 for one ARKitScenes validation seed. A potential explanation is that the previous loss overweights appearance texture, but this has not been tested causally. SACG without boundary loss still uses edge consistency inside its gate.

The simple CNN is a meaningful and competitive control. It narrowly trails SACG in L1 and AbsRel but improves RMSE, RGB-edge BoundaryAbsRel and SSIM relative to the selected method. Because the CNN consumes metric LiDAR values directly whereas SACG refinement uses frozen DPT features, this is not a clean isolated test of gating.

Limitations include no calibrated uncertainty supervision, no explicit preservation of valid LiDAR observations in refined depth, no guarantee of positive output depth, no independent scene-held-out test, no three-seed results for the selected objective, and no selected-model ScanNet++ evaluation. The evaluation split also influenced method choice.

In summary, SACG combines frozen PromptDA with a learned validity-guided spatial gate and a DPT-feature residual. The no-boundary training variant shows promising but small single-seed improvements. Robust claims about architectural advantage require additional controlled ablations and independent evaluation.
\bibliographystyle{plain}
\bibliography{references}
\end{document}
